\begin{pastexam}[label={px:ncku-112-prog-04}]{成大資訊聯招 112 學年度 第 4 題}
（複選）

{[5\%]} Given an AVL tree as follows, we perform the operations (A), (B), (C), (D) and (E),
independently. Which of these operations would result in the tree to be unbalanced, and shall be
restructured?

\begin{center}
\begin{tikzpicture}[x=1cm,y=1cm,>={Stealth[length=2mm]},every node/.style={circle,draw,inner sep=0pt,minimum size=7mm,font=\sffamily\small}]
\node (a) at (0,0) {22};
\node (b) at (-0.98,-1.1) {17};
\node (c) at (0.84,-1.1) {35};
\node (d) at (-0.29,-2.37) {20};
\node (e) at (1.53,-2.37) {36};
\draw[->] (a)--(b); \draw[->] (a)--(c); \draw[->] (b)--(d); \draw[->] (c)--(e);
\end{tikzpicture}
\end{center}
\begin{enumerate}[label=(\Alph*)]
\item insert 18
\item insert 33
\item insert 38
\item delete 35
\item delete 22
\end{enumerate}
\end{pastexam}

\begin{pastsoln}
答案：(A)(C)。

原樹：根 $22$；左子 $17$，$17$ 只有右子 $20$；右子 $35$，$35$ 只有右子 $36$。 高度用節點數算（資結書
第~7.1~節），各節點的平衡因子（左高減右高）：$20$、$36$ 是 $0$，$17$、$35$ 是 $-1$，$22$ 是 $0$。 每個操作
都從原樹獨立做，做完從變動處往上看有沒有平衡因子變成 $\pm 2$（資結書第~10.2~節）。

\begin{itemize}
\item (A) 插入 $18$：$18 < 22$、$> 17$、$< 20$，成為 $20$ 的左子。 $17$ 的左子樹高 $0$、右子樹（$20, 18$）
高 $2$，平衡因子 $-2$，失衡，要調整（RL 型：先在 $20$ 右旋，再在 $17$ 左旋，$18$ 升上來）。 是。
\item (B) 插入 $33$：成為 $35$ 的左子，$35$ 變成 $0$，$22$ 的右子樹高度不變。 不失衡。
\item (C) 插入 $38$：成為 $36$ 的右子。 $35$ 的左子樹高 $0$、右子樹（$36, 38$）高 $2$，平衡因子 $-2$，失衡
（RR 型，在 $35$ 左旋）。 是。
\item (D) 刪除 $35$：$35$ 只有一個子節點 $36$，$36$ 接上來。 $22$ 的左子樹高 $2$、右子樹高 $1$，平衡因子
$+1$。 不失衡。
\item (E) 刪除 $22$：兩個子節點，用中序後繼 $35$ 取代（$35$ 的位置由 $36$ 接上），根 $35$ 的左子樹高 $2$、
右子樹高 $1$，平衡。 若改用中序前驅 $20$ 取代，根 $20$ 的左子樹（$17$）高 $1$、右子樹（$35, 36$）高 $2$，
也平衡。 不失衡。
\end{itemize}
\end{pastsoln}
